\section{The smallest Hasse-principle failure}\label{sec:hasse}

For a number field \(K\), call a Keller fiber \(A\) a \emph{Hasse failure} if
\(\Spec A\) has a point over every completion of \(K\) but no \(K\)-rational
point.  In this section \(d_{\mathrm{HP}}\) denotes the least rank of such a
fiber over \(\Q\).  We show that \(d_{\mathrm{HP}}=5\).

\begin{lemma}[Degree-four barrier]\label{lem:degree-four-barrier}
A nonempty finite \'etale $\Q$-scheme of degree at most four which has a point
over every $\Q_p$ has a rational point.
\end{lemma}

\begin{proof}
Let $X$ be such a scheme and let $L$ split it.  At every unramified prime,
$X(\Q_p)\neq\varnothing$ exactly when the Frobenius element fixes a point of
the finite $\operatorname{Gal}(L/\Q)$-set $X(L)$.

Suppose first that $X$ is connected of degree greater than one.  The Galois
action is transitive and contains a derangement.  Indeed, Burnside's lemma says
that the average number of fixed points is one, whereas the identity fixes
more than one point.  Some group element therefore fixes none.  Chebotarev
produces an unramified prime with that Frobenius class, contradicting local
solubility.

Now suppose that $X$ is disconnected and has no rational component.  In total
degree at most four, the only remaining possibility is a disjoint union of two
quadratic fields.  If they are equal, choose an inert prime.  If they are
distinct, their compositum has Galois group $C_2\times C_2$ and contains an
element nontrivial in both quadratic quotients.  Chebotarev again gives a prime
inert in both components.  Thus every everywhere locally soluble finite
\'etale scheme of degree at most four contains $\Spec\Q$.
\end{proof}

We use the classical polynomial of Berend and Bilu,
\begin{equation}
 P_5(T)=(T^3-19)(T^2+T+1).
 \label{eq:BB}
\end{equation}
It appeared in their study of polynomials with roots modulo every integer
\cite{BerendBilu}.

\begin{theorem}[Optimal Hasse-failing Keller fiber]\label{thm:hasse}
There is an explicit Keller map $F:\A^3_\Q\to\A^3_\Q$ of geometric degree
five and a rational target $y$ such that
\[
 F^{-1}(y)(\Q)=\varnothing,
 \qquad
 F^{-1}(y)(\Q_v)\neq\varnothing
 \quad\text{for every place }v.
\]
No Hasse-failing Keller fiber over \(\Q\) has rank below five.  Hence
\[
 \boxed{d_{\mathrm{HP}}=5}.
\]
\end{theorem}

\begin{proof}
Write
\[
 g(T)=T^3-19,\qquad q(T)=T^2+T+1.
\]
The following table is the complete nonarchimedean local-solubility audit.
In a generic row, the displayed residue root is simple and hence lifts by
Hensel's lemma.  For a finite base-case audit we remove \(2,3,19\) from the
generic rows and give a direct witness at each of them.
\begin{center}
\small
\begin{tabular}{c|c|c|c}
prime $p$ & factor & residue or approximate root & Hensel check\\
\hline
$p\notin\{2,3,19\}$, $p\equiv2\pmod3$
 & $g$ & $a^3=19$ in $\mathbb F_p$
 & $g'(a)=3a^2\neq0$\\
$p\notin\{2,3,19\}$, $p\equiv1\pmod3$
 & $q$ & $\zeta\in\mathbb F_p^\times$, $\operatorname{ord}(\zeta)=3$
 & $q'(\zeta)=2\zeta+1\neq0$\\
$p=2$ & $g$ & $a=1$ & $g(1)=-18$, $g'(1)=3\in\mathbb Z_2^\times$\\
$p=3$ & $g$ & $a=-2$ &
 $v_3(g(-2))=3>2=2v_3(g'(-2))$\\
$p=19$ & $q$ & $a=7$ &
 $q(7)=57$, $q'(7)=15\in\mathbb Z_{19}^\times$
\end{tabular}
\end{center}
Here are the justifications behind the two generic rows.  If
$p\equiv2\pmod3$, then $\gcd(3,p-1)=1$, so cubing is an automorphism of
$\mathbb F_p^\times$; since $p\neq19$, it gives a nonzero $a$ with
$a^3=19$, and $p\neq3$ makes $3a^2$ nonzero.  If $p\equiv1\pmod3$, the
cyclic group $\mathbb F_p^\times$ contains an element $\zeta$ of order three.
Then $\zeta\neq1$ and
\[
 0=\zeta^3-1=(\zeta-1)q(\zeta).
\]
The root is simple because the two roots of $q$ can coincide only when its
discriminant $-3$ vanishes, which is excluded.  The rows $p=2$ and $p=19$
use ordinary simple-root Hensel.  The $p=3$ row uses the strong form of
Hensel's lemma: $g(-2)=-27$ and $g'(-2)=12$, so
$v_3(g(-2))>2v_3(g'(-2))$.

This covers every finite prime, since every prime outside
$\{2,3,19\}$ is congruent to $1$ or $2$ modulo $3$.  At the real place,
$g(2)=-11<0<8=g(3)$, so the intermediate value theorem supplies a real
root.

There is no rational root: the rational-root theorem excludes a root of
$g$ (none of $\pm1,\pm19$ is a root), while $q$ has discriminant $-3$.
Moreover, $P_5=gq$ is squarefree: both factors are squarefree and they are
coprime, since a root of $q$ has cube $1$ and therefore cannot be a root of
$g$.  The map below has geometric degree five and, by
Corollary~\ref{cor:squarefree-gauge}, has fiber exactly
$\Spec\Q[T]/(P_5)$.  Its local and global points are consequently exactly
those just audited.  Finally, Lemma~\ref{lem:degree-four-barrier} applies to
every complete regular fiber of rank at most four and proves rank-minimality.
\end{proof}

For reference, here is the actual map.  It comes directly from the normalized
gauge, so its determinant does not require a separate guess or search.  Put
\[
 G(T)=P_5(T)-P_5(0)=T^5+T^4+T^3-19T^2-19T.
\]
Thus
\[
 (g_1,g_2,g_3,g_4,g_5)=(-19,-19,1,1,1).
\]
With $t=1+xy$ and
\[
 q=t^2z-19y^2(1+3t),
\]
Theorem~\ref{thm:localized-gauge} gives the gauge map
\begin{align*}
 \Pi={}&tq,\\
 B_0={}&y-\frac3{19}xq+2tq
       -\frac4{19}t^2x^2q^4-\frac5{19}t^2x^3q^5,\\
 C_0={}&x(5-3t)+\frac1{19}x^3z
       +\frac2{19}(xq)^4+\frac3{19}(xq)^5,
\end{align*}
with $\det D(\Pi,B_0,C_0)=-2$.  Apply the target linear automorphism
\[
 L(\Pi,B_0,C_0)=(\Pi,19B_0,19C_0).
\]
The resulting denominator-free map $F=L\circ F_G$ is
\begin{equation}
\begin{aligned}
 F_1={}&tq,\\
 F_2={}&19y-3xq+38tq-4t^2x^2q^4-5t^2x^3q^5,\\
 F_3={}&19x(5-3t)+x^3z+2(xq)^4+3(xq)^5.
\end{aligned}
\label{eq:explicit-hasse-map}
\end{equation}
Consequently,
\[
 \boxed{\det DF=\det(L)(-2)=19^2(-2)=-722.}
\]
The normalized target $(1,0,-2)$ is sent to
\[
 y=(1,0,-38).
\]
At that target the inverse equation is
\[
 G(T)-\frac{-19}{2}(-2)=G(T)-19=P_5(T).
\]
The finite-\'etale (squarefree) fiber corollary,
Corollary~\ref{cor:squarefree-gauge}, therefore gives
\begin{equation}
 \boxed{F^{-1}(1,0,-38)\simeq
 \Spec\Q[T]/\bigl((T^3-19)(T^2+T+1)\bigr).}
 \label{eq:explicit-hasse-fiber}
\end{equation}
So the smallest Hasse-failing Keller fiber is not only existential; its map,
target, determinant, and quotient algebra are all explicit.
